\section*{第 \ref{ch:g9:acids-bases-ph} 章——酸、碱与 pH}

\begin{solution}{exo:g9:acids-bases-ph:1}
pH 3：酸性。pH 7：中性。pH 11：碱性。
\end{solution}

\begin{solution}{exo:g9:acids-bases-ph:2}
酸性：\ce{H+} 离子。碱性：\ce{OH-} 离子。
\end{solution}

\begin{solution}{exo:g9:acids-bases-ph:3}
青柠汁（2）、咖啡（5）、纯水（7）、海水（8.1）、
家用氨水（12）。
\end{solution}

\begin{solution}{exo:g9:acids-bases-ph:4}
取一条试纸放在干净干燥的表面皿上，
用干净玻璃棒取一滴溶液点在试纸上，
立即与比色卡比较颜色。
\end{solution}

\begin{solution}{exo:g9:acids-bases-ph:5}
稀释十倍后约为 5；稀释一百倍后约为 6。
\end{solution}

\begin{solution}{exo:g9:acids-bases-ph:6}
不能。稀释会使 pH 更接近 7，但永远不会越过 7：
水本身呈中性，加水不可能使 \ce{OH-} 离子
的数目多于 \ce{H+} 离子。
\end{solution}

\begin{solution}{exo:g9:acids-bases-ph:7}
混合时会放出热量。若把水倒在浓酸上，
水可能立即沸腾，并把酸溅出；
而把酸慢慢倒入水中，酸则会立即分散到大量水中。
\end{solution}

\begin{solution}{exo:g9:acids-bases-ph:8}
腐蚀性：会灼伤皮肤和眼睛，并腐蚀金属。
应佩戴护目镜和手套；溅出的液体要用大量水冲洗。
\end{solution}

\begin{solution}{exo:g9:acids-bases-ph:9}
三次稀释（1 至 2、2 至 3、3 至 4）：$10 \times 10 \times 10 =
1000$ 倍。
\end{solution}

\begin{solution}{exo:g9:acids-bases-ph:10}
氢氧化镁乳剂呈碱性（pH 约为 10.5）：
它能与胃中过量的部分酸发生反应。
\end{solution}

\begin{solution}{exo:g9:acids-bases-ph:11}
不能：pH 为 8.1 时它仍呈弱碱性。它正在
向 pH 更低的方向移动、碱性减弱，
这是因为空气中的二氧化碳溶入其中。
\end{solution}

\begin{solution}{exo:g9:acids-bases-ph:12}
pH 试纸需用肉眼对照比色卡读取，精确到约一个单位；
pH 计则可给出一位或两位小数。pH 计更精确；
两者在试纸的精度范围内是一致的。
\end{solution}

\begin{solution}{pb:g9:acids-bases-ph:1}
\textbf{1.} 酸性：其中主要含有 \ce{H+(aq)} 和 \ce{Cl-(aq)}
离子。

\textbf{2.} GHS05，腐蚀性：酸会腐蚀皮肤、眼睛和金属。

\textbf{3.} 它给出精确的数值，
而不是供肉眼对照的颜色。

\textbf{4.} 多十倍：pH 每相差一个单位，即对应十倍稀释。

\textbf{5.} \qty{10}{L}，pH 为 3。

\textbf{6.} 三次：pH 从 2 到 3、从 3 到 4、从 4 到 5。

\textbf{7.} $1 \times 10 \times 10 \times 10 = \qty{1000}{L}$。

\textbf{8.} 不能：稀释只会使 pH 更接近 7。

\textbf{9.} \ce{H+ + OH- -> H2O}。

\textbf{10.} 酸被消耗并转化为水，而不是
分散到极大的体积中；只需少量碱，
无需上千升的水。

\textbf{11.} 溶于水的小苏打（或肥皂水）。

\textbf{12.} $1000 - 1 = \qty{999}{L}$ 的水。
\end{solution}
